\documentclass[aspectratio=169, 10pt]{beamer}
\usepackage{beamer_theme_cienciadados2026}

\title[Aula 03 - Paradigmas e Tipos]{Ciência dos Dados I: Análise Exploratória}
\subtitle{Aula 03: Paradigmas de Programação e Sistemas de Tipagem}
\author{Prof. Dr. Ney Lemke}
\institute{UNESP -- Instituto de Biociências de Botucatu}
\date{Versão 2026}

\begin{document}

\frame{\titlepage}

\begin{frame}{Sumário da Aula}
  \tableofcontents
\end{frame}

\section{O que é um Paradigma de Programação?}

\begin{frame}{Classificação dos Paradigmas}
  Um paradigma representa uma forma fundamental de estruturar o raciocínio computacional:
  \begin{itemize}
    \item \textbf{Imperativo / Procedural}: Foco no \textit{como fazer} (passo a passo de instruções e alteração de estado na memória). Ex: C, Fortran, Pascal.
    \item \textbf{Funcional}: Foco no \textit{o que calcular} (composição de funções puras, ausência de efeitos colaterais e dados imutáveis). Ex: Haskell, Lisp, Elixir.
    \item \textbf{Orientado a Objetos (OOP)}: Agrupamento de dados (atributos) e comportamentos (métodos) em entidades coesas (classes e objetos). Ex: Java, C++, Smalltalk.
  \end{itemize}
  \begin{block}{Python é Multiparadigma}
    Python adota pragmatismo: permite combinar o melhor de cada paradigma conforme a necessidade da análise de dados!
  \end{block}
\end{frame}

\section{Sistema de Tipagem em Python}

\begin{frame}{Dinâmica, Forte e Type Hints}
  \begin{columns}
    \begin{column}{0.5\textwidth}
      \begin{block}{Tipagem Dinâmica}
        O tipo da variável é inferido em tempo de execução. Uma mesma variável pode referenciar diferentes tipos.
      \end{block}
      \begin{block}{Tipagem Forte}
        Python não realiza coerção implícita de tipos incompatíveis (ex: \texttt{'10' + 5} resulta em \texttt{TypeError}).
      \end{block}
    \end{column}
    \begin{column}{0.48\textwidth}
      \begin{block}{Type Hints (PEP 484 até 2026)}
        Anotações explícitas de tipo aumentam a legibilidade e evitam bugs em pipelines complexos de dados.
      \end{block}
    \end{column}
  \end{columns}
\end{frame}

\section{Exemplo Comparativo: Imperativo vs. Funcional}

\begin{frame}[fragile]{Soma dos Quadrados dos Pares}
  \begin{columns}
    \begin{column}{0.5\textwidth}
      \textbf{Estilo Imperativo (Loops e Mutação)}
      \begin{lstlisting}
soma = 0
for x in range(1, 11):
    if x % 2 == 0:
        soma += x ** 2
print(soma)
      \end{lstlisting}
    \end{column}
    \begin{column}{0.5\textwidth}
      \textbf{Estilo Funcional (Composição)}
      \begin{lstlisting}
pares = filter(lambda x: x % 2 == 0, range(1, 11))
quadrados = map(lambda x: x ** 2, pares)
soma = sum(quadrados)
print(soma)
      \end{lstlisting}
    \end{column}
  \end{columns}
  \begin{exampleblock}{Na Prática de Dados}
    Pandas e PySpark utilizam princípios funcionais intensamente (transformações encadeadas sem alteração dos dados originais).
  \end{exampleblock}
\end{frame}

\begin{frame}{Resumo da Aula}
  \begin{itemize}
    \item Paradigmas moldam como expressamos transformações de dados.
    \item Tipagem forte garante integridade e type hints documentam expectativas.
    \item \textbf{Prática}: Expressões lógicas, casting explícito e funções lambda.
  \end{itemize}
\end{frame}

\end{document}
